\begin{frame}{PPO in Practice}
\begin{itemize}
    \item The clipped objective is only half the story --- these implementation details matter \emph{as much as} $\varepsilon$, and explain why two ``correct'' PPO implementations can differ by an order of magnitude\footnotemark
    \vspace{0.3em}
    \item \textbf{Advantage normalisation:} standardise $\hat{A}_t$ per minibatch (stabilises gradient scale)
    \item \textbf{Value-loss clipping:} clip the critic update like the policy
    \item \textbf{Gradient clipping:} bound global grad-norm (typically $0.5$)
    \item \textbf{KL early stopping:} end the epoch loop once mean KL exceeds a target --- a cheap trust-region safeguard
    \item \textbf{Obs / reward normalisation:} running normalisation keeps inputs and returns well-scaled
    \item \textbf{Orthogonal init:} small policy-head gain avoids an over-confident initial policy
\end{itemize}
\footnotetext{Andrychowicz et al.\ (2020), \emph{What Matters in On-Policy RL?}}
\end{frame}
